\chapter{Problem Statement and Research Questions}
\label{ch:rqs}

The background chapter established two things: architectural documentation drifts from implementation, and the implementation carries several distinct signals that could substitute for that documentation, typed dependencies, folder organisation, lexical identifiers, and member-level detail, none of which is sufficient on its own. This chapter turns that observation into the research questions that organise the rest of the thesis. The questions are stated at the level of the thesis as a whole rather than at the level of an individual paper; each paper in Part I and Part II answers a narrower, internally scoped question, and the mapping between the two is given in Section~\ref{sec:orq-mapping}.

Two axes structure the work. The first axis concerns \emph{signal and representation}: given that architectural information is distributed across several implementation artefacts, which signals carry that information, how should they be represented for a graph-based learner, and what limits each one. The second axis concerns \emph{task}: whether a finding obtained under unsupervised architecture recovery, where no part of the architecture is known in advance, also holds under semi-supervised code-to-architecture mapping, where a small labelled subset anchors the rest. The five studies in this thesis populate a two-by-something grid along these axes rather than five independent investigations, and the overarching research questions below are written to make that grid explicit.

\section{Overarching research questions}
\label{sec:orqs}

\paragraph{ORQ1: Do learned graph representations recover software architecture more effectively than hand-tuned multi-signal fusion, and where does the advantage concentrate?}
Existing architecture recovery techniques that combine multiple signals, such as SARIF, do so through manually weighted similarity functions tuned on the systems at hand. This question asks whether replacing that hand-tuned fusion with a learned encoder, trained to reconstruct the observed dependency structure while conditioning on semantic and organisational node features, yields better alignment with ground-truth architectures, and under what system properties the improvement is largest or smallest.

\paragraph{ORQ2: What does each implementation signal, typed dependencies, folder organisation, lexical identifiers, and member-level detail, contribute to recovery, and what bounds each signal's usefulness?}
This is the thesis's central question, and it is deliberately broad because no single study answers it alone. Typed dependencies carry structural coupling but are bounded by how cohesively the ground-truth architecture is reflected in the dependency graph, and by whether the model's neighbourhood reach matches the target granularity. Folder organisation reflects how developers already partition responsibility and is often a strong proxy for the intended decomposition, but adopting it shifts recovery away from dependency-based modularity, and it stops being useful once its own hierarchy is too fine or too flat to disambiguate membership. Lexical identifiers encode developer intent, but the studies in this thesis find that conventional representations of this signal, bag-of-words and averaged Word2Vec embeddings, carry comparatively little architectural information on their own, and that representation quality, not just the presence of the signal, determines how much it can contribute. Member-level detail refines file-level representations when file-level dependencies alone are too coarse. The chapters in Parts I and II each isolate a subset of these signals under controlled conditions; the discussion in Chapter~\ref{ch:discussion} draws the contribution and the boundary condition for each together.

\paragraph{ORQ3: Is the choice of encoder or the model's access to structural context the binding constraint on dependency-based recovery?}
A recurring pattern across the unsupervised studies is that a shallow encoder applied directly to the raw dependency graph underperforms not because the encoder is the wrong architecture, but because two-hop message passing is too local to expose the longer-range structure that separates architectural modules. This question asks whether expanding a model's effective reach, through random-walk context, diffusion-based rewiring, or additional message-passing depth, closes the gap attributed to encoder choice, and whether this pattern generalises from unsupervised recovery to supervised mapping.

\paragraph{ORQ4: Do the answers to ORQ1 through ORQ3 hold across both unsupervised architecture recovery and semi-supervised code-to-architecture mapping?}
The thesis evaluates overlapping signals, graph constructions, and benchmark systems under two different tasks and two different evaluation regimes: alignment to a ground-truth partition in the unsupervised setting, and macro/micro F1 against partially known labels in the supervised setting. This question asks whether the same properties, encoder choice being secondary to reach, folder organisation compensating for weak dependency cohesion, lexical representation quality mattering more than the presence of lexical features, transfer between the two tasks, or whether the supervised setting surfaces behaviour the unsupervised one does not, such as the graph's effect on decision repeatability rather than on raw accuracy.

\paragraph{ORQ5: How strongly does directory structure alone stand in for the intended architecture, and what can compensate for it when it is unavailable?}
This question is subordinate to ORQ2 and is stated separately only because two studies address it directly from complementary angles: one by isolating folder organisation from dependency structure without lexical features, and the other by removing folder organisation entirely from a richer heterogeneous graph and asking what representation of the remaining signals recovers the lost ground. It is not, on its own, the thesis's main claim, and Chapter~\ref{ch:discussion} treats it as one instance of ORQ2 rather than as an independent finding.

\section{Mapping research questions to studies}
\label{sec:orq-mapping}

Table~\ref{tab:orq-mapping} summarises which studies feed which overarching question. A study can feed more than one question, and a question is answered cumulatively rather than by any single study.

% Table sketch (fill once numbering of papers/chapters is final):
% ORQ1 <- ARL, GAER
% ORQ2 <- all five, with StructSAR and HetMap as the primary signal-isolation studies
% ORQ3 <- StructSAR primarily, supported by the C2A extension's depth ablation and HetMap's encoder choice
% ORQ4 <- the pairing of Part I and Part II as a whole; most directly the C2A extension against GAER/StructSAR
% ORQ5 <- StructSAR (folder without lexical features) and HetMap (folder removed from the full graph)

\section{Methodological approach common to the studies}
\label{sec:common-method}

The five studies share a methodological spine that is stated once here rather than repeated in each chapter. Dependencies are extracted with the multilingual static-analysis tool Depends, aggregated to file or class granularity depending on the study. Node features combine, in varying subsets depending on the research question under test, folder-path tokens, Word2Vec embeddings of normalised code identifiers, and, in the case of HetMap, embeddings obtained through dependency-guided contrastive fine-tuning. Benchmark systems are drawn from a common pool of open-source Java, C, and C++ projects with published ground-truth code-to-architecture mappings, ranging from small systems of a few hundred files to Chromium at roughly eighteen thousand files, so that findings can be checked against system size, dependency density, and ground-truth granularity rather than against a single representative case. Because the learning methods are stochastic, each configuration is run repeatedly, typically 100 or 200 times with different random seeds, and results are reported as medians with run-to-run variability shown through interquartile ranges or boxplots. Pairwise comparisons use nonparametric tests, the Wilcoxon signed-rank or Mann-Whitney U test depending on whether the compared samples are paired, together with Cliff's delta to report effect size, since the sample sizes involved are large enough that statistical significance alone would be uninformative. Finally, the number of recovered clusters, or the target module count in the supervised setting, is treated throughout as a modelling choice that materially affects every metric, rather than as a nuisance parameter to be fixed once and ignored.